% --- Section: Abstract ---

\begin{abstract}
Failure modes are disguised on accuracy-only leaderboards, and the
disguise matters when LLMs are used to solve statistical problems.
Five frontier models (Claude Sonnet 4.5, GPT-4.1, Gemini 2.5 Flash,
DeepSeek-Chat, Mistral Large) are benchmarked against 171 Bayesian
and inferential tasks across 38 types with 473 perturbations on
three axes: rephrase, semantic, and numerical. Each response is
scored twice, by a deterministic keyword rubric and by an
out-of-family judge (Llama 3.3 70B Turbo, served via Together AI)
on assumption-compliance and reasoning quality. The rankings
disagree. GPT-4.1 places last on accuracy (0.621) but first on
robustness (0.9996). Claude leads accuracy (0.683) and binary-pass
calibration (ECE 0.0334), yet drops to fourth on robustness
(0.9565). 44.3\% of responses do not mention the assumptions the
derivation depends on, and assumption violation drives 47\% of
judge-classified failures. The keyword rubric and the judge agree
only moderately on assumption-compliance (Spearman 0.602) and
barely at all on the two other judged dimensions. A single-number
leaderboard is the wrong instrument here: which model wins depends
on the metric.
\end{abstract}

\begin{IEEEkeywords}
LLMs, Bayes' theorem, calibration, perturbation robustness,
expected calibration error, LLM-as-judge, benchmark
\end{IEEEkeywords}
